% Source: Wald GR, physical PDF pp. 325--337 (printed pp. 312--324).
% Coverage: Section 12.3, The Charged Kerr Black Holes; equations (12.3.1)--(12.3.30).
% Figures/tables/footnotes: Figures 12.3--12.6; footnotes 3--5; no tables.
% Status: source-reviewed and compile-clean.
% Uncertainties: none.

\subsection{The Charged Kerr Black Holes}\label{sec:12.3}
\setcounter{theorem}{0}

Consider a body which undergoes complete gravitational collapse and---in accord with
the first cosmic censor conjecture of Section~12.1---forms a black hole. In the
spherically symmetric case, the spacetime outside the body always is described by the
Schwarzschild solution, and the final state of gravitational collapse will be a
Schwarzschild black hole. However, in the nonspherical case, the spacetime geometry
outside the collapsing body should vary with time and depend greatly on the details of
the collapse. In particular, no gravitational radiation can be produced in a
spherically symmetric spacetime, whereas large amounts of energy may be radiated
away in nonspherical collapse. Nevertheless, one would expect on physical grounds
that at sufficiently \lq\lq late times,\rq\rq\ the spacetime geometry should
\lq\lq settle down\rq\rq\ to a stationary final state. Furthermore, one would expect
all the matter present to be rapidly \lq\lq swallowed up\rq\rq\ by the black hole, so
the final state should be vacuum except, possibly, for the presence of electromagnetic
fields associated with the black hole. (Even in cases---such as in the binary X-ray
sources---where one has a steady flow of matter into the black hole, this matter
normally would produce only a small perturbation on the structure of the black hole.)
Thus, one expects that the final state of gravitational collapse will be a stationary,
electrovac (i.e., vacuum except for electromagnetic fields) black hole. Therefore, it
is of great interest to find all solutions of the Einstein--Maxwell equations which
describe stationary black holes.

We have already discussed in detail the spherically symmetric, static black hole
solution discovered by Schwarzschild (1916a). Shortly thereafter, a charged
generalization of the Schwarzschild solution was discovered independently by
Reissner (1916) and Nordstrom (1918) (see problem 3 of chapter 6). However, it was
not until 1963 that another family of stationary, vacuum black hole solutions was
discovered by Kerr. A charged generalization of the Kerr family was obtained shortly
thereafter by Newman et al.\ (1965). These charged Kerr solutions form a
three-parameter family, whose spacetime metric and electromagnetic vector potential
are given by
\begin{equation}
\begin{aligned}
 \dd s^2={}&-\left(\frac{\Delta-a^2\sin^2\theta}{\Sigma}\right)\dd t^2
 -\frac{2a\sin^2\theta(r^2+a^2-\Delta)}{\Sigma}\dd t\dd\phi\\
 &+\left[\frac{(r^2+a^2)^2-\Delta a^2\sin^2\theta}{\Sigma}\right]
 \sin^2\theta\,\dd\phi^2+\frac{\Sigma}{\Delta}\dd r^2+\Sigma\dd\theta^2 ,
\end{aligned}
\tag{12.3.1}\label{eq:12.3.1}
\end{equation}
\begin{equation}
 A_a=-\frac{er}{\Sigma}\bigl[(\dd t)_a-a\sin^2\theta(\dd\phi)_a\bigr],
 \tag{12.3.2}\label{eq:12.3.2}
\end{equation}
where
\begin{align}
 \Sigma&=r^2+a^2\cos^2\theta, \tag{12.3.3}\label{eq:12.3.3}\\
 \Delta&=r^2+a^2+e^2-2Mr. \tag{12.3.4}\label{eq:12.3.4}
\end{align}
Here \(e\), \(a\), and \(M\) are the three parameters of the family. When \(e=0\),
we have \(A_a=0\) and the spacetime metric reduces to the vacuum Kerr family of
solutions. When \(a=0\), we recover the Reissner--Nordstrom solutions, and when
\(e=a=0\), equation~\eqref{eq:12.3.1} reduces to the Schwarzschild solution. Thus,
all known stationary black hole solutions are encompassed by the three-parameter
family \eqref{eq:12.3.1}, \eqref{eq:12.3.2}. As we shall see at the end of this
section, no other stationary black hole solutions exist.

The charged Kerr metrics all are stationary and axisymmetric (see Section~7.1), with
Killing fields \(\xi^a=(\partial/\partial t)^a\) and
\(\psi^a=(\partial/\partial\phi)^a\). They are asymptotically flat, as can be seen
crudely from the fact that the metric components \eqref{eq:12.3.1} approach those
of the Minkowski metric in spherical polar coordinates as \(r\to\infty\), and has
been demonstrated in detail by Ashtekar and Hansen (1978). In the algebraic
classification of Section~7.3, they are type II--II solutions, with repeated principal
null vectors,
\begin{align}
 l^a&=\frac{r^2+a^2}{\Delta}\left(\frac{\partial}{\partial t}\right)^a+
      \frac a\Delta\left(\frac{\partial}{\partial\phi}\right)^a+
      \left(\frac{\partial}{\partial r}\right)^a, \tag{12.3.5}\label{eq:12.3.5}\\
 n^a&=\frac{r^2+a^2}{2\Sigma}\left(\frac{\partial}{\partial t}\right)^a+
      \frac a{2\Sigma}\left(\frac{\partial}{\partial\phi}\right)^a-
      \frac\Delta{2\Sigma}\left(\frac{\partial}{\partial r}\right)^a .
      \tag{12.3.6}\label{eq:12.3.6}
\end{align}
(Here \(l^a\) and \(n^a\) are normalized by a convenient choice due to Kinnersley
1969, with \(l^an_a=-1\).)

The three parameters \(e\), \(a\), and \(M\) appearing in the solutions all have a
direct physical interpretation. For any 2-sphere, \(S\), in the asymptotic region, we
have
\begin{equation}
 \frac12\int_S\epsilon_{abcd}F^{cd}=4\pi e ,
 \tag{12.3.7}\label{eq:12.3.7}
\end{equation}
so, by problem 2 of chapter 4, we may interpret \(e\) as the total electric charge of
the spacetime. Furthermore, we have
\begin{equation}
 -\frac1{8\pi}\int_S\epsilon_{abcd}\nabla^c\xi^d=M ,
 \tag{12.3.8}\label{eq:12.3.8}
\end{equation}
so, according to equation~(11.2.9), \(M\) is the total mass. Finally, we have
\begin{equation}
 \frac1{16\pi}\int_S\epsilon_{abcd}\nabla^c\psi^d=Ma ,
 \tag{12.3.9}\label{eq:12.3.9}
\end{equation}
so, according to problem 6 of chapter 11, we have \(a=J/M\), where \(J\) is the
total angular momentum of the spacetime.

It should be noted that, in geometrized units, the charge-to-mass ratio of a proton is
\(q/m\sim10^{18}\), and for an electron we have \(q/m\sim10^{21}\). Since the ratio
of electromagnetic to gravitational force produced on a test body of charge \(q\) and
mass \(m\) by a body of charge \(e\) and mass \(M\) is \(\sim qe/mM\), it would be
very difficult for any astrophysical body to achieve and/or maintain a charge-to-mass
ratio of greater than \(\sim10^{-18}\), since a body with larger charge-to-mass ratio
would selectively attract particles of opposite charge.\footnote{One mechanism to
obtain a net charge on an astrophysical body which, in principle, could overcome this
limit is to have a rotating body placed in a magnetic field. However, for a Kerr black
hole, this mechanism also would lead to a negligible charge buildup in astrophysically
reasonable situations (Wald 1974b).} Hence, in astrophysically reasonable situations
it appears that \(e\ll M\), so we may neglect the effects of the electromagnetic field
on the spacetime geometry and consider only the Kerr family of black holes.

The coordinate basis components of the charged Kerr metrics,
equation~\eqref{eq:12.3.1}, are nonsingular and define a nondegenerate metric
everywhere except where \(\Sigma=0\) and where \(\Delta=0\). Evaluation of
curvature invariants such as \(R_{abcd}R^{abcd}\) shows that the singularity at
\begin{equation}
 \Sigma=r^2+a^2\cos^2\theta=0
 \tag{12.3.10}\label{eq:12.3.10}
\end{equation}
is a true, curvature singularity when \(M\ne0\). If one were to interpret \(r\),
\(\theta\), and \(\phi\) as representing spherical polar coordinates, the fact that
there is a singularity at the origin, \(r=0\), only for the angular value
\(\theta=\pi/2\) would appear rather puzzling. In particular, if we interpret the
singularity in this way---i.e., if we define the charged Kerr metrics on the manifold
\(\mathbb R^4\) with the origin \(r=0\) removed---we then would have incomplete
geodesics (such as those on the axis, \(\sin\theta=0\)) which terminate at \(r=0\)
but along which the curvature remains finite. In fact, this spacetime would be
extendible. This provides a good illustration of the impropriety of making a choice of
the manifold structure on the basis of a naive interpretation of the coordinate system
in which the metric is given.

Some insight into the true nature of the singularity of the charged Kerr metrics at
\(\Sigma=0\) can be obtained by consideration of the case \(e=M=0\), \(a\ne0\).
In that case, the Kerr metric \eqref{eq:12.3.1} actually is nothing more than the
metric of Minkowski spacetime expressed in spheroidal coordinates. Here the
singularity at \(\Sigma=0\) is, of course, merely a coordinate singularity, and it is
located on the ring of radius \(a\) in the plane \(z=0\). This suggests that when
\(M\ne0\) and \(a\ne0\) we interpret the true singularity at \(\Sigma=0\) as a ring
singularity, i.e., that we define the charged Kerr metrics on a manifold whose
structure in a neighborhood of this singularity has the topology of \(\mathbb R^4\)
with the set \(S^1\times\mathbb R\)---that is, a ring, \(S^1\), cross
\lq\lq time,\rq\rq\ \(\mathbb R\)---removed. This can be implemented explicitly by
transforming to the quasi-Cartesian coordinates given by Kerr and Schild (1965),
where \(\Sigma=0\) takes the coordinate form of a ring. However, a problem still
remains in that when \(M\ne0\) the metric components fail to be smooth across the
coordinate disk enclosed by the ring singularity in the \(z=0\) plane. This problem
can be remedied by defining the charged Kerr metrics on a manifold with the
following relatively complicated topology in a neighborhood of the singularity at
\(\Sigma=0\). We take two copies, \(M_1\) and \(M_2\), of \(\mathbb R^4\) with the
\lq\lq ring\rq\rq\ \(z=0\), \(x^2+y^2=a^2\), removed. We then attach \(M_1\) to
\(M_2\) by identifying the \lq\lq top side\rq\rq\ of the disk
\(z=0\), \(x^2+y^2<a^2\), of \(M_1\) with the \lq\lq bottom side\rq\rq\ of the
corresponding disk of \(M_2\), and, similarly, identify the \lq\lq bottom side\rq\rq\
of the disk of \(M_1\) with the \lq\lq top side\rq\rq\ of the disk of \(M_2\). The
charged Kerr metrics with \(M\ne0\), \(a\ne0\) then may be smoothly defined on
this manifold in such a way that the curvature scalar \(R_{abcd}R^{abcd}\) blows up
along every incomplete geodesic, thereby guaranteeing that the spacetime is
inextendible. Details of this construction can be found in Hawking and Ellis (1973).

It should be noted that when one passes \lq\lq through the ring\rq\rq\ in going from
\(M_1\) to \(M_2\) in this spacetime, this corresponds in the original coordinates to
passing through \(r=0\) into negative values of \(r\). However, for negative values
of \(r\) of sufficiently small magnitude and for \(\theta\) sufficiently close to
\(\pi/2\) we have \(\psi_a\psi^a=g_{\phi\phi}<0\) (see
equation~\eqref{eq:12.3.1}). Thus, \(\psi^a=(\partial/\partial\phi)^a\) becomes
timelike near the ring singularity. However, the orbits of \(\psi^a\) must be closed
(i.e., the coordinate \(\phi\) must be periodically identified with period \(2\pi\)) in
order that the charged Kerr spacetime be asymptotically flat as \(r\to\infty\). Thus,
closed timelike curves exist in a neighborhood of the ring singularity.

When \(e^2+a^2>M^2\), there are no solutions of the equation \(\Delta=0\) so the
true singularity at \(\Sigma=0\) is the only singularity of the coordinate components
\eqref{eq:12.3.1}. In this case, the ring singularity is \lq\lq naked,\rq\rq\ i.e., the
charged Kerr metrics fail to be strongly asymptotically predictable, and thus they do
\emph{not} describe black holes. Furthermore, one may make use of the causality
violation occurring near the ring singularity to go \lq\lq backwards in time\rq\rq\ by
an arbitrarily large amount as measured by the \(t\) coordinate of \eqref{eq:12.3.1}
and thereby produce closed timelike curves passing through any point in the
spacetime. In the case
\begin{equation}
 e^2+a^2\le M^2
 \tag{12.3.11}\label{eq:12.3.11}
\end{equation}
\(\Delta\) vanishes at the \(r\)-coordinate values
\begin{equation}
 r_\pm=M\pm(M^2-a^2-e^2)^{1/2}.
 \tag{12.3.12}\label{eq:12.3.12}
\end{equation}

As shown by Boyer and Lindquist (1967) and Carter (1968a), the singularities in the
metric components at \(r=r_+\) and (for \(a\ne0\) or \(e\ne0\)) at \(r=r_-\) are
coordinate singularities of the same nature as the singularity at \(r=2M\) in the
Schwarzschild spacetime. Thus, we may extend the spacetime through these coordinate
singularities much as in the Schwarzschild case. When these extensions are patched
together, a remarkable global structure of the extended charged Kerr spacetimes is
obtained. A conformal diagram of the extended Schwarzschild spacetime is shown in
Figure~12.3, and a conformal diagram of the extended charged Kerr spacetime with
\(a\ne0\) is shown in Figure~12.4 for the \lq\lq non-extreme\rq\rq\ case
\(a^2+e^2<M^2\). Region I of Figure~12.4 is the asymptotically flat region covered
in a nonsingular fashion by the original coordinates \eqref{eq:12.3.1} with
\(r>r_+\). By extending through the coordinate singularity at \(r=r_+\), we obtain
region II representing a black hole, region III representing a white hole, and region
IV representing another asymptotically flat region, just as in the Schwarzschild case,
Figures~6.9 and~12.3. However, unlike the Schwarzschild case, instead of
encountering a true singularity at the \lq\lq top boundary\rq\rq\ of region II and
\lq\lq bottom boundary\rq\rq\ of region III, we encounter merely another coordinate
singularity at \(r=r_-\). Thus, we can extend region II through \(r=r_-\) to obtain
regions V and VI. These regions contain the ring singularity at \(\Sigma=0\) and, as
described above, one can pass through the ring singularity to obtain another
asymptotically flat region with \(r\to-\infty\). (In this asymptotically flat region
the ring singularity is a naked singularity of negative mass. With respect to the
original asymptotically flat region I, the ring singularity, of course, lies within a
black hole.) One may then continue to extend the charged Kerr spacetime
\lq\lq upward\rq\rq\ ad infinitum to obtain a region VII, identical in structure to
region III, and obtain regions VIII and IX, identical in structure to regions IV and I,
etc. Similarly, one may extend the charged Kerr solutions \lq\lq downward\rq\rq\ ad
infinitum. The structure of the extended Reissner--Nordstrom spacetime
(\(a=0,e\ne0\)) is very similar except that the true singularity at \(\Sigma=0\) no
longer has a ring structure, and one cannot extend to negative values of \(r\). The
global structure of the \lq\lq extreme\rq\rq\ charged Kerr case \(e^2+a^2=M^2\)
(where \(r_+=r_-=M\)) differs from Figure~12.4 but has a similar structure
consisting of \lq\lq blocks\rq\rq\ with \(r>M\) and \(r<M\) patched together in an
infinite chain.

% Source: Wald GR, physical PDF p. 329
%         (printed p. 316).
% Coverage: Conformal diagram of the extended Schwarzschild spacetime.
% Figures/tables/footnotes: Figure 12.3; no tables or footnotes.
% Status: source-reviewed and compile-clean.
% Uncertainties: none.

\begingroup
\renewcommand{\thefigure}{12.3}
\begin{figure}[H]
  \centering
  \begin{tikzpicture}[x=.80cm,y=.80cm,>=Latex,line cap=round,line join=round,font=\small]
    \coordinate (C) at (0,0);
    % Four adjoining Schwarzschild regions.
    \draw[thick] (C) -- (-1.45,1.45) -- (-2.90,0) -- (-1.45,-1.45) -- cycle;
    \draw[thick] (C) -- (1.45,1.45) -- (2.90,0) -- (1.45,-1.45) -- cycle;
    \draw[thick] (C) -- (-1.45,1.45) -- (0,2.90) -- (1.45,1.45);
    \draw[thick] (C) -- (-1.45,-1.45) -- (0,-2.90) -- (1.45,-1.45);
    \draw[thick,decorate,decoration={zigzag,segment length=2.1pt,amplitude=.7pt}]
      (-1.45,1.45) -- (0,2.90) -- (1.45,1.45);
    \draw[thick,decorate,decoration={zigzag,segment length=2.1pt,amplitude=.7pt}]
      (-1.45,-1.45) -- (0,-2.90) -- (1.45,-1.45);
    \node at (1.53,.13) {I};
    \node at (-1.53,.13) {IV};
    \node at (0,1.92) {II};
    \node at (0,-1.92) {III};
    \node at (2.55,.96) {\(\mathcal I^+\)};
    \node at (2.55,-.96) {\(\mathcal I^-\)};
    \node at (-2.55,.96) {\(\mathcal I^+\)};
    \node at (-2.55,-.96) {\(\mathcal I^-\)};
    \node at (3.22,0) {\(\ii^0\)};
    \node at (-3.22,0) {\(\ii^0\)};
  \end{tikzpicture}
  \caption{A conformal diagram of the extended Schwarzschild spacetime (see
  Figure~6.9), represented in the same manner as used in Figure~12.2. Note that
  since the extended Schwarzschild spacetime has two distinct asymptotically flat
  regions, two distinct conformal boundaries are shown.}
  \label{fig:12.3}
\end{figure}
\endgroup

% Source: Wald GR, physical PDF p. 330
%         (printed p. 317).
% Coverage: Conformal diagram of the extended charged Kerr spacetime.
% Figures/tables/footnotes: Figure 12.4; no tables or footnotes.
% Status: source-reviewed and compile-clean.
% Uncertainties: none.

\begingroup
\renewcommand{\thefigure}{12.4}
\begin{figure}[H]
  \centering
  \begin{tikzpicture}[x=1.50cm,y=1.15cm,>=Latex,line cap=round,line join=round,font=\small]
    % The repeated vertical Carter chain.  Adjacent diamonds meet at single
    % vertices; they are not three independent horizontal stacks.
    \newcommand{\diamondat}[2]{%
      \draw[thick] (#1-.85,#2) -- (#1,#2+.85) -- (#1+.85,#2) -- (#1,#2-.85) -- cycle;}
    \diamondat{-0.85}{-1.70}
    \diamondat{ 0.85}{-1.70}
    \diamondat{ 0.00}{-0.85}
    \diamondat{-0.85}{ 0.00}
    \diamondat{ 0.85}{ 0.00}
    \diamondat{ 0.00}{ 0.85}
    \diamondat{-0.85}{ 1.70}
    \diamondat{ 0.85}{ 1.70}
    \diamondat{ 0.00}{ 2.55}
    \diamondat{-0.85}{ 3.40}
    \diamondat{ 0.85}{ 3.40}
    \node at (-1.16,0) {IV};      \node at (1.16,0) {I};
    \node at (0,.85) {II};        \node at (0,-.85) {III};
    \node at (-1.18,1.72) {\(V'\)}; \node at (-.52,1.72) {V};
    \node at ( .52,1.72) {VI};      \node at (1.18,1.72) {\(VI'\)};
    \node at (0,2.55) {VII};
    \node at (-.85,3.40) {VIII};  \node at (.85,3.40) {IX};
    % Representative horizon and asymptotic labels, placed on the same edges as
    % in the source without obscuring the causal regions.
    \node[rotate=45,fill=white,inner sep=.4pt] at (-1.33,.44) {\tiny \(r=\infty\)};
    \node[rotate=-45,fill=white,inner sep=.4pt] at (1.33,.44) {\tiny \(r=\infty\)};
    \node[rotate=-45,fill=white,inner sep=.4pt] at (-1.33,-.44) {\tiny \(r=\infty\)};
    \node[rotate=45,fill=white,inner sep=.4pt] at (1.33,-.44) {\tiny \(r=\infty\)};
    \node[rotate=45,fill=white,inner sep=.4pt] at (-.53,.30) {\tiny \(r=r_+\)};
    \node[rotate=-45,fill=white,inner sep=.4pt] at (.53,.30) {\tiny \(r=r_+\)};
    \node[rotate=-45,fill=white,inner sep=.4pt] at (-.53,-.30) {\tiny \(r=r_+\)};
    \node[rotate=45,fill=white,inner sep=.4pt] at (.53,-.30) {\tiny \(r=r_+\)};
    \node[rotate=-45,fill=white,inner sep=.4pt] at (-.53,1.40) {\tiny \(r=r_-\)};
    \node[rotate=45,fill=white,inner sep=.4pt] at (.53,1.40) {\tiny \(r=r_-\)};
    \node[rotate=45,fill=white,inner sep=.4pt] at (-1.33,2.14) {\tiny \(r=-\infty\)};
    \node[rotate=-45,fill=white,inner sep=.4pt] at (1.33,2.14) {\tiny \(r=-\infty\)};
    % Initial data surface and ring singularity.
    \draw[very thick] (-1.68,-.02) .. controls (-.85,.55) and (.85,.55) .. (1.68,-.02);
    \node[left] at (-1.70,-.35) {\(S\)};
    \draw[dashed] (-.85,1.37) -- (-.85,1.95);
    \draw[dashed] ( .85,1.37) -- ( .85,1.95);
    \draw[->] (2.12,1.40) .. controls (1.72,1.48) and (1.34,1.47) .. (.92,1.44);
    \node[right,align=left] at (2.16,1.16) {\(\Sigma=0\)\\(Ring Singularity)};
  \end{tikzpicture}
  \caption{A conformal diagram of the extended charged Kerr spacetime in the
  case \(a\ne0\), \(a^2+e^2<M^2\).}
  \label{fig:12.4}
\end{figure}
\endgroup


Thus, an observer starting in region I of the extended charged Kerr spacetime of
Figure~12.4 may cross the event horizon at \(r=r_+\) and enter the black hole region
II. However, instead of inevitably falling into a singularity within a finite proper time
as occurs in the Schwarzschild spacetime, the observer may pass through the
\lq\lq inner horizon\rq\rq\ \(r=r_-\) (which is a Cauchy horizon for the hypersurface
\(S\) shown in Fig.~12.4), thereby entering region V or VI. From there, he may end
his existence in the ring singularity but he also may pass through the ring singularity
to a new asymptotically flat region, or he may enter the \lq\lq white hole region\rq\rq\
VII and from there enter the new asymptotically flat region VIII or IX. From there he
may enter the new black hole associated with these asymptotic regions, and continue
his journey.

How much of this extended charged Kerr spacetime should be taken seriously? What
portion of this spacetime would be produced by a physically realistic gravitational
collapse, starting from \lq\lq non-exotic\rq\rq\ initial data, i.e., from an
asymptotically flat initial data surface \(S\) with topology \(\mathbb R^3\)? Unlike
the Schwarzschild case, we have no reason to believe that the exterior gravitational
field of any physically reasonable collapsing body will be described by the Kerr
metric, since, as mentioned above, in the nonspherical case we would, in general,
expect a complicated dynamical evolution which only \lq\lq settles down\rq\rq\ to a
stationary geometry at late times in \(J^-(\mathcal I^+)\). Thus, we are not in a
position to follow the dynamical evolution of the gravitational collapse of a body
which forms a Kerr black hole and thereby determine the detailed spacetime geometry
inside the black hole. However, in the case of spherical collapse of a charged body
(\(e\ne0\)), the spacetime geometry exterior to the matter \emph{is} described by the
Reissner--Nordstrom solution since Birkhoff's theorem can be generalized to show that
the Reissner--Nordstrom spacetime is the unique spherical electrovac solution. The
dynamical evolution of a simple system like a charged, spherical shell of dust can be
obtained explicitly. In this case, spacetime is flat inside the shell, and the flat
interior of the shell entirely \lq\lq covers up\rq\rq\ regions III and IV of
Figure~12.4. Part or all of regions II and V (including, in all cases, the singularity
at \(r=0\) in region V) also are covered up. The behavior of the shell for the various
choices of total mass \(M\), total charge \(e\), and total rest mass \(\mathcal M\) are
chronicled in detail by Boulware (1973). We show the resulting spacetime for the
case \(M>\mathcal M>|e|\) in Figure~12.5.

% Source: Wald GR, physical PDF p. 331
%         (printed p. 318).
% Coverage: Conformal diagram of collapse of a charged spherical shell.
% Figures/tables/footnotes: Figure 12.5; no tables or footnotes.
% Status: source-reviewed and compile-clean.
% Uncertainties: none.

\begingroup
\renewcommand{\thefigure}{12.5}
\begin{figure}[H]
  \centering
  \begin{tikzpicture}[x=.76cm,y=.76cm,>=Latex,line cap=round,line join=round,font=\small]
    \coordinate (B) at (0,-2.35); \coordinate (P) at (-.18,1.20);
    \coordinate (R) at (2.15,0);
    \draw[thick] (B) -- (R) -- (P);
    \node at (.95,-.16) {I};
    \node at (-.02,.72) {II};
    \node at (-.42,1.30) {VI};
    % The collapsing shell and extended Reissner--Nordstrom continuation.
    \draw[very thick] (-.38,2.15) .. controls (-1.45,.72) and (-1.12,-.92) .. (B);
    \draw[very thick] (-.68,2.27) .. controls (-.42,.58) and (-.70,-.62) .. (B);
    \draw[dashed] (-2.65,.65) -- (-.80,-1.10) -- (-2.65,-2.15);
    \draw[dashed] (-2.10,1.52) -- (-.53,.07) -- (-2.25,-.95);
    \draw[thick,decorate,decoration={zigzag,segment length=2.1pt,amplitude=.7pt}]
      (-.38,2.15) -- (.12,2.72);
    \node[above right] at (.02,2.58) {\(r=0\) (Singularity)};
    \draw[->] (-3.0,1.80) .. controls (-2.0,1.70) and (-1.35,1.47) .. (-.82,1.17);
    \node[left,align=right] at (-3.02,2.04) {Surface\\of Shell};
    \draw[->] (-2.55,.45) .. controls (-1.85,.35) and (-1.30,.25) .. (-.86,.12);
    \node[left,align=right] at (-2.58,.53) {\(r=0\)\\(Origin of Coordinates)};
    \node[right] at (.55,.88) {\(r=r_-\)};
    \node[right] at (.78,.20) {\(r=r_+\)};
    \node[right] at (2.55,.90) {\(\mathcal I^+\)};
    \node[right] at (2.55,-.90) {\(\mathcal I^-\)};
    \node[right] at (2.35,0) {\(i^0\)};
  \end{tikzpicture}
  \caption{A conformal diagram of a spacetime in which a charged spherical shell
  with \(M>\mathcal M>|e|\) undergoes gravitational collapse. The dashed lines
  refer to the extended Reissner--Nordstrom spacetime.}
  \label{fig:12.5}
\end{figure}
\endgroup


As shown there, the shell crosses the surface \(r=r_-\), which is a Cauchy horizon
for region I. The spacetime is extendible across \(r_-\), but the extension is not
determined by Einstein's equation since it is outside the domain of dependence of the
initial data surface. However, if one assumes that the extension is spherically
symmetric, the extension is uniquely given by the Reissner--Nordstrom geometry, and
we find that the shell crashes into the singularity at \(r=0\) (which formed
\emph{outside} the shell) in region VI, as shown in Figure~12.5. For other choices of
the parameters \(M\), \(\mathcal M\), and \(e\), the shell may reexpand into region VII
of the Reissner--Nordstrom geometry (Boulware 1973). Thus, one may obtain
spacetimes from the collapse of a charged body which have features differing greatly
from that of Figure~6.11. However, as already mentioned in our discussion of cosmic
censorship in Section~12.1, it has been shown that the Cauchy horizon of
Figure~12.5 is unstable (Chandrasekhar and Hartle 1982). Linear perturbations of the
initial data for the Einstein--Maxwell equations on a Cauchy surface for region I of
Figure~12.5 become singular at \(r=r_-\). The basic reason for this is that an observer
crossing \(r=r_-\) in Figure~12.5 \lq\lq sees\rq\rq\ all of region I. Oscillations in
the gravitational and electromagnetic field which occur at finite frequency in region I
are \lq\lq seen\rq\rq\ to occur at infinite frequency by an observer crossing \(r=r_-\);
i.e., there is an \lq\lq infinite blueshift\rq\rq\ effect, which makes the perturbation
singular there. Thus, there is good reason to believe that in a physically realistic case
where the shell is not exactly spherical, the Cauchy horizon \(r=r_-\) in
Figure~12.5 will become a true, physical singularity, thereby producing an
\lq\lq all encompassing\rq\rq\ singularity inside the black hole formed by the collapse
of the shell. It is believed that similar phenomena will occur for any physically
realistic collapse to a charged Kerr black hole. Thus, the spacetime produced by
physically realistic collapse is expected to be qualitatively similar to the spherical
case, Figure~6.11, rather than that suggested by the extended charged Kerr solutions,
Figure~12.4.

Another feature of the charged Kerr spacetime worthy of note is that the norm of the
timelike Killing field,
\begin{equation}
 \xi^a\xi_a=g_{tt}=\frac{a^2\sin^2\theta-\Delta}{\Sigma},
 \tag{12.3.13}\label{eq:12.3.13}
\end{equation}
becomes positive in the region where
\begin{equation}
 r^2+a^2\cos^2\theta+e^2-2Mr<0,
 \tag{12.3.14}\label{eq:12.3.14}
\end{equation}
part of which lies \emph{outside} the black hole if \(a\ne0\). Thus, in the region
\begin{equation}
 r_+<r<M+(M^2-e^2-a^2\cos^2\theta)^{1/2},
 \tag{12.3.15}\label{eq:12.3.15}
\end{equation}
called the \emph{ergosphere} and depicted in Figure~12.6, the asymptotic time
translation Killing field \(\xi^a=(\partial/\partial t)^a\) becomes spacelike. Hence,
an observer in the ergosphere would have to \lq\lq go faster than light\rq\rq\ to
follow an orbit of \(\xi^a\); i.e., he cannot remain stationary, even though he is
outside the black hole. The nature of this nonstationarity can be seen from the
equation
\begin{equation}
 \sum_{\mu,\nu}g_{\mu\nu}u^\mu u^\nu<0
 \tag{12.3.16}\label{eq:12.3.16}
\end{equation}
satisfied by the components of the tangent vector, \(u^a\), to any timelike curve.
Inside the ergosphere, all terms on the left-hand side of \eqref{eq:12.3.16} are
manifestly positive except the term
\(2g_{t\phi}u^tu^\phi=2g_{t\phi}(\dd t/\dd\tau)(\dd\phi/\dd\tau)\), where
\(\tau\) denotes proper time along the curve. Since \(\nabla^at\) is past directed
timelike in the ergosphere, we have \(\dd t/\dd\tau=u^a\nabla_at>0\). Thus, since
\(g_{t\phi}<0\) in the ergosphere, we must have
\begin{equation}
 \dd\phi/\dd t>0
 \tag{12.3.17}\label{eq:12.3.17}
\end{equation}
for all timelike curves in the ergosphere. In other words, all observers in the
ergosphere are forced to rotate in the direction of rotation of the black hole. This may
be viewed as an extreme case of the \lq\lq dragging of inertial frames\rq\rq\ effect,
thereby providing a dramatic example of how some aspects of Mach's principle are
incorporated into general relativity.

The closest analog to a family of static observers outside the black hole in the charged
Kerr geometry are the \lq\lq locally nonrotating observers\rq\rq\ (see problem 3 of
chapter 7) whose 4-velocity is given by
\(u^a=-\nabla^at/[-\nabla_bt\nabla^bt]^{1/2}\). These observers rotate with
coordinate angular velocity
\begin{equation}
 \Omega=\frac{\dd\phi}{\dd t}=-\frac{g_{t\phi}}{g_{\phi\phi}}
 =\frac{a(r^2+a^2-\Delta)}{(r^2+a^2)^2-\Delta a^2\sin^2\theta}.
 \tag{12.3.18}\label{eq:12.3.18}
\end{equation}

% Source: Wald GR, physical PDF p. 332
%         (printed p. 319).
% Coverage: Side and top views of a Kerr ergosphere.
% Figures/tables/footnotes: Figure 12.6; no tables or footnotes.
% Status: source-reviewed and compile-clean.
% Uncertainties: none.

\begingroup
\renewcommand{\thefigure}{12.6}
\begin{figure}[H]
  \centering
  \begin{tikzpicture}[x=.75cm,y=.75cm,>=Latex,line cap=round,line join=round,font=\small]
    % Side view.
    \draw[thick] (-4.0,0) ellipse [x radius=1.72,y radius=.67];
    \draw[thick] (-4.0,0) circle [radius=.72];
    \draw[thick] (-4.0,.20) -- (-4.0,1.70);
    \draw[->,thick] (-4.28,1.45) arc[start angle=200,end angle=520,x radius=.28,y radius=.11];
    \draw[->] (-2.58,.18) .. controls (-3.05,.16) and (-3.15,.06) .. (-3.34,.00);
    \node[right] at (-2.47,.44) {Black Hole};
    \draw[->] (-2.55,-1.05) .. controls (-3.05,-.82) and (-3.22,-.50) .. (-3.30,-.33);
    \node[right] at (-2.42,-1.08) {Ergosphere};
    \node at (-4.0,-1.85) {(a)};
    % Top view.
    \draw[thick] (3.0,0) circle [radius=1.37];
    \draw[thick] (3.0,0) circle [radius=.72];
    \draw[->,thick] (2.55,.15) arc[start angle=170,end angle=35,radius=.46];
    \node at (3.0,-1.85) {(b)};
  \end{tikzpicture}
  \caption{A sketch showing (a) a \lq\lq side view\rq\rq\ and (b) a \lq\lq top view\rq\rq\ of
  the ergosphere of a Kerr black hole.}
  \label{fig:12.6}
\end{figure}
\endgroup


In the limit as one approaches the black hole event horizon, \(r\to r_+\), this
coordinate angular velocity becomes
\begin{equation}
 \Omega_H=\frac{a}{r_+^2+a^2}.
 \tag{12.3.19}\label{eq:12.3.19}
\end{equation}
This is closely related to the fact that it is the Killing field
\begin{equation}
 X^a=\left(\frac{\partial}{\partial t}\right)^a+
       \Omega_H\left(\frac{\partial}{\partial\phi}\right)^a
 \tag{12.3.20}\label{eq:12.3.20}
\end{equation}
[rather than \((\partial/\partial t)^a\)] which is tangent to the null geodesic
generators of the horizon of the charged Kerr black hole. Equation~\eqref{eq:12.3.20}
can be interpreted as saying that the event horizon of the charged Kerr black hole
rotates with angular velocity \(\Omega_H\).

We turn, now, to a brief discussion of geodesic motion. For simplicity, we shall
consider only the Kerr geometry, \(e=0\). (See problem 2 for the case \(e\ne0\).) As
in the Schwarzschild case, the timelike Killing field \(\xi^a\) and the axial Killing
field \(\psi^a\) yield via proposition~C.3.1 a conserved energy, \(E\), and angular
momentum, \(L\), per unit rest mass for geodesics,
\begin{align}
 E&=-u^a\xi_a=\left(1-\frac{2Mr}{\Sigma}\right)\dot t+
       \frac{2Mar\sin^2\theta}{\Sigma}\dot\phi,
       \tag{12.3.21}\label{eq:12.3.21}\\
 L&=u^a\psi_a=-\frac{2Mar\sin^2\theta}{\Sigma}\dot t+
       \frac{(r^2+a^2)^2-\Delta a^2\sin^2\theta}{\Sigma}
       \sin^2\theta\,\dot\phi,
       \tag{12.3.22}\label{eq:12.3.22}
\end{align}
where \(\dot x^\mu=\dd x^\mu/\dd\tau\). In addition, we have
\begin{equation}
 g_{ab}u^au^b=-\kappa,
 \tag{12.3.23}\label{eq:12.3.23}
\end{equation}
where \(\kappa=1\) for timelike geodesics and \(\kappa=0\) for null geodesics. One
may use equations~\eqref{eq:12.3.21} and~\eqref{eq:12.3.22} to eliminate \(\dot t\)
and \(\dot\phi\) in terms of \(E\) and \(L\), and the result may be substituted into
equation~\eqref{eq:12.3.23}. In the case of equatorial geodesics,
\(\theta=\pi/2\), one obtains
\begin{equation}
 \frac12\dot r^2+V(E,L,r)=0,
 \tag{12.3.24}\label{eq:12.3.24}
\end{equation}
where
\begin{equation}
 V=-\kappa\frac Mr+\frac{L^2}{2r^2}
  +\frac12(\kappa-E^2)\left(1+\frac{a^2}{r^2}\right)
  -\frac M{r^3}(L-aE)^2 .
 \tag{12.3.25}\label{eq:12.3.25}
\end{equation}
Thus, as in the Schwarzschild case, the problem of obtaining the timelike and null
geodesics in the equatorial plane of the Kerr spacetime reduces to solving a problem
of ordinary nonrelativistic, one-dimensional motion in an effective potential, with
the only significant additional complication here being the fact that \(V\) now
depends nontrivially on \(E\) as well as on \(L\). Thus, we may find the behavior of
freely falling test bodies and light rays by methods similar to those of Section~6.3.
In particular, the circular orbits are given by the simultaneous solutions of \(V=0\)
and \(\dd V/\dd r=0\). Their properties are discussed by Bardeen, Press, and
Teukolsky (1972). It is noteworthy that binding energies significantly higher than in
the Schwarzschild case can be achieved for circular orbits around a Kerr black hole.
For a Kerr black hole with \(a=M\), the last stable circular orbit (with positive \(L\))
has \(E=1/\sqrt3\). Thus, if a test body with positive \(L\) spirals in to the last
stable circular orbit as a result of energy loss via gravitational radiation, it will have
radiated away \(1-1/\sqrt3\simeq42\%\) of its original rest energy, as opposed to
only \(\sim6\%\) in the Schwarzschild case, equation~(6.3.23).

The constants of motion \(E\) and \(L\), equations~\eqref{eq:12.3.21} and
\eqref{eq:12.3.22}, do \emph{not} provide a sufficient number of first integrals to
determine nonequatorial motion. Therefore, in that case we might expect to have to
return to the geodesic equation, \(u^a\nabla_au^b=0\), and solve a coupled set of
second order nonlinear ordinary differential equations to obtain \(r(\tau)\) and
\(\theta(\tau)\). Remarkably, however, it turns out that this is not necessary. The
Kerr metric possesses the Killing tensor,
\begin{equation}
 K_{ab}=2\Sigma l_{(a}n_{b)}+r^2g_{ab},
 \tag{12.3.26}\label{eq:12.3.26}
\end{equation}
(Walker and Penrose 1970), where \(l^a\) and \(n^a\) are given by
equations~\eqref{eq:12.3.5} and~\eqref{eq:12.3.6}. Thus, an additional constant of
the motion
\begin{equation}
 C=K_{ab}u^au^b
 \tag{12.3.27}\label{eq:12.3.27}
\end{equation}
is obtained, as discussed at the end of appendix C. This allows one to integrate the
geodesic equation explicitly, as was first done by Carter (1968a), who used the
separability of the Hamilton--Jacobi equation for geodesics rather than the existence
of \(K_{ab}\) to obtain this additional constant of motion. In addition, the Kerr
metric---as well as all other type II--II vacuum spacetimes---possesses a
\lq\lq conformal Killing spinor\rq\rq\ (see Walker and Penrose 1970) which enables
one to determine the parallel propagation of \lq\lq polarization vectors\rq\rq\ along
null geodesics in a simple manner.

A remarkable simplification also occurs when one studies the propagation of a
Klein--Gordon scalar test field, equation~(4.3.9), in the Kerr spacetimes. One can
take advantage of the stationary and axial Killing fields of Kerr to expand the scalar
field in a Fourier series in the angular coordinate \(\phi\) and a Fourier integral over
the time coordinate \(t\). This effectively reduces the Klein--Gordon equation to a
partial differential equation in the two remaining variables, \(r\) and \(\theta\).
However, it turns out that this equation can be solved by separation of variables
(Carter 1968b), so the problem of determining the behavior of Klein--Gordon test
fields in the Kerr spacetimes is reduced to solving ordinary differential equations.

Even more remarkable simplifications occur when one studies Maxwell's equations
and the linearized Einstein equation in the Kerr background. Here, after making use
of the symmetries in \(\phi\) and \(t\), one would expect to be left with a complicated,
coupled system of partial differential equations in \(r\) and \(\theta\) for,
respectively, the components of the electromagnetic vector potential, \(A_\mu\), and
the components of the metric perturbation, \(h_{\mu\nu}\). Indeed, this is what
happens when one writes down the equations for \(A_\mu\) or \(h_{\mu\nu}\), even
when simplifying gauge choices such as the Lorenz gauge for \(A_\mu\) or the
transverse traceless gauge for \(h_{\mu\nu}\) (see Section~7.5) are made. However,
one may write down the equations in the Newman--Penrose (1962) formalism (see
Section~3.4b), using the repeated principal null vectors of the Kerr metric \(l^a\)
and \(n^a\), equations~\eqref{eq:12.3.5} and~\eqref{eq:12.3.6}, as the real null
vectors and
\begin{equation}
 m^a=\frac{1}{2^{1/2}(r+\ii a\cos\theta)}
 \left[\ii a\sin\theta\left(\frac{\partial}{\partial t}\right)^a+
       \left(\frac{\partial}{\partial\theta}\right)^a+
       \frac{\ii}{\sin\theta}\left(\frac{\partial}{\partial\phi}\right)^a\right]
 \tag{12.3.28}\label{eq:12.3.28}
\end{equation}
as the complex null vector of the Newman--Penrose basis. Then, as discovered by
Teukolsky (1972), a decoupled equation can be derived for
\begin{equation}
 \Phi_0=F_{ab}l^am^b
 \tag{12.3.29}\label{eq:12.3.29}
\end{equation}
in the Maxwell case, and for
\begin{equation}
 \Psi_0=-C_{abcd}l^am^bl^cm^d
 \tag{12.3.30}\label{eq:12.3.30}
\end{equation}
in the linearized Einstein case. (Decoupled equations similarly may be obtained for
\(\Phi_2=F_{ab}\bar m^an^b\) and
\(\Psi_4=-C_{abcd}n^a\bar m^bn^c\bar m^d\).) Furthermore, Teukolsky showed that
these equations also may be solved by separation of variables. In addition, a
knowledge of \(\Phi_0\) (or \(\Phi_2\)) determines a Maxwell perturbation modulo
the trivial \lq\lq change in charge\rq\rq\ solution obtained by linearizing
equation~\eqref{eq:12.3.2} off the Kerr background, while a knowledge of \(\Psi_0\)
(or \(\Psi_4\)) determines a gravitational perturbation modulo the trivial
perturbations obtained by variation of the Kerr parameters \(M,a\) (Wald 1973).
Finally, the complete vector potential, \(A_a\), and metric perturbation, \(h_{ab}\),
associated, respectively, with \(\Phi_0\) and \(\Psi_0\), can be obtained explicitly
(Cohen and Kegeles 1974; Wald 1978a; Chandrasekhar 1983). Thus, the problem of
determining the behavior of electromagnetic and gravitational perturbations of a Kerr
black hole also reduces to solving ordinary differential equations, thereby making
tractable many problems of interest. Similar simplifications of the coupled linearized
Einstein--Maxwell system occur in the Reissner--Nordstrom case, \(a=0,e\ne0\)
(Moncrief 1975; Chandrasekhar 1979). However, it appears that no such
simplifications occur in the general charged Kerr case, \(a\ne0,e\ne0\). For a
complete discussion of the electromagnetic and gravitational perturbations of a Kerr
black hole, we refer the reader to Chandrasekhar (1983).

Thus, as summarized above, a great deal is known about the properties of the Kerr
black holes. These are the only known stationary vacuum black hole solutions of
Einstein's equation. However, since they comprise only a two-parameter family, one
might expect that many more stationary vacuum black hole solutions should exist.
After all, as mentioned in the introduction to chapter 11, the exterior gravitational
field of a stationary body is characterized by an infinite set of multipole coefficients.
Why should all the higher multipole moments of a stationary black hole be related in a
unique way to its mass and angular momentum? Remarkably, as a result of the
theorems of Israel, Carter, Hawking, and Robinson obtained between 1967 and 1975,
a virtually complete proof has been given that the Kerr black holes are the only
possible stationary vacuum black holes. Thus, if the first cosmic censor conjecture is
correct and if the spacetime resulting from gravitational collapse always
\lq\lq settles down\rq\rq\ to a stationary, vacuum final state, the end product of
collapse must always be a Kerr black hole. The complete gravitational collapse of two
bodies differing greatly from each other in composition, shape, and structure will
produce indistinguishable final states provided only that their end products have the
same total mass and total angular momentum.

Since we are far from having the general stationary vacuum solution of Einstein's
equation in explicit form,\footnote{Indeed, until the early 1970s the Kerr solutions
were virtually the only known stationary, nonstatic, asymptotically flat vacuum
solutions. As discussed in Sections~7.1 and~7.4, great progress has been made in
obtaining the general stationary, axisymmetric vacuum solution, but even so we are
far from having the solutions in sufficiently explicit form to determine if they
represent black holes.} the proof of the uniqueness of the Kerr black holes has
proceeded by a relatively long chain of arguments. Logically, the first step---although,
historically, one of the last steps---in this chain is the proof by Hawking that the
two-dimensional surface formed by the intersection of the horizon of a stationary
black hole with a Cauchy surface must have topology \(S^2\). This is established by
showing that if it had any other topology, it would be possible to deform it outward
into \(J^-(\mathcal I^+)\) such that the expansion, \(\theta\), of the outgoing null
geodesics satisfies \(\theta\le0\) everywhere. This would contradict
proposition~12.2.4. Details of the proof can be found in Hawking and Ellis (1973).
(See also Gannon 1976 for some results in the nonstationary case.)

The next step in the uniqueness proof, also due to Hawking, is the demonstration that
a stationary vacuum black hole must be static or axisymmetric. First, we note that in
a stationary spacetime containing a black hole, the time translation isometry must
leave the horizon invariant. Hence, the Killing field \(\xi^a\) must lie tangent to the
horizon and, hence, always must be spacelike or null on the horizon. Now, one of the
following three possibilities must hold: (i) No ergosphere is present in the spacetime,
i.e., the stationary Killing field \(\xi^a\) is everywhere timelike or null outside the
black hole. In this case, \(\xi^a\) must be null on the horizon. (ii) An ergosphere is
present but is disjoint from the horizon, and \(\xi^a\) is null on the horizon. (iii) An
ergosphere is present and intersects the horizon, as happens for a Kerr black hole.
In this case, \(\xi^a\) is spacelike on (a portion of) the horizon. In case (i), a
generalization of a theorem of Lichnerowicz (1955) establishes that the spacetime
must be static (Hawking and Ellis 1973). Under some additional assumptions,
results of Hajicek (1973) show that the outer boundary of the ergosphere in a
stationary vacuum spacetime always must intersect the horizon. Thus, it appears that
case (ii) cannot occur. Plausibility arguments against case (ii) also are given in
Hawking and Ellis (1973). Finally, in case (iii), using the properties of the horizon in
a stationary spacetime and using the analyticity of stationary vacuum solutions
(Müller zum Hagen 1970), Hawking proved existence of a one-parameter group of
isometries which commute with the stationary isometries and whose orbits on the
horizon coincide with the null geodesic generators of the horizon. Thus, one obtains a
Killing field \(X^a\) distinct from \(\xi^a\), and by taking a linear combination of
\(X^a\) and \(\xi^a\), one obtains a Killing field \(\psi^a\) whose orbits are closed,
i.e., an axial Killing field.\footnote{The proof that a stationary, nonstatic black hole
must be axisymmetric continues to hold in the case where a distribution of matter is
placed outside a rotating black hole. This leads to an apparent paradox since one would
expect it to be possible to \lq\lq hold in place\rq\rq\ a nonaxisymmetric distribution
of matter far from the black hole, thereby producing a stationary nonaxisymmetric
spacetime. The resolution of this paradox is that such a matter distribution will
produce an effective \lq\lq tidal friction\rq\rq\ causing the black hole to
\lq\lq spin down\rq\rq\ and thus be nonstationary until it reaches a final static state.
A discussion of this process is given by Hawking and Hartle (1972).} Again, details
of this proof are given in Hawking and Ellis (1973).

The case of a static, vacuum, topologically spherical black hole was analyzed by
Israel (1967), who proved that the only such black holes are the Schwarzschild
solutions. Some additional assumptions were made in the proof, but the most notable
of them---that all the surfaces of constant \(\xi^a\xi_a\) are topologically
spheres---has been eliminated by Müller zum Hagen, Robinson, and Seifert (1973)
and Robinson (1977).

Finally the case of a stationary, axisymmetric, vacuum topologically spherical black
hole was analyzed by Carter (1971) and Robinson (1975) using the methods described
in Section~7.1 to cast Einstein's equation and the black hole boundary conditions into
a relatively simple form. They succeeded in proving that all stationary axisymmetric
black holes are uniquely characterized by two parameters which appear in the boundary
conditions. Since the Kerr solutions exhaust all possible values of these parameters,
it follows that they are the only possible stationary axisymmetric black holes.

The above results have been generalized to the electrovac case. Hawking's theorem
on the spherical topology of a stationary black hole still applies since it requires only
that the dominant energy condition be satisfied by matter. The proof that a stationary
black hole with no ergosphere must be static generalizes to the electrovac case
(Carter 1973), and the proof of existence of an axial Killing field if \(\xi^a\) is
spacelike on the horizon depends only on the general form, (10.1.21), of the equations
and thus also applies to the electrovac case. Israel's theorem has been generalized to
show that the only possible static, electrovac black holes are the Reissner--Nordstrom
solutions (Israel 1968). Finally, Mazur (1982) and Bunting (unpublished) have
generalized the Carter--Robinson proof of uniqueness of Kerr to show that the charged
Kerr solutions (together with their generalizations possessing magnetic charge, which
are obtained by applying a duality rotation to the charged Kerr electromagnetic field)
are the only stationary, axisymmetric electrovac solutions. Further generalizations to
exclude the possible presence of other types of classical fields around a black hole
also have been given (Bekenstein 1972; Hartle 1972; Teitelboim 1972). In addition,
numerous examples have been given (see, e.g., Wald 1972a) to illustrate how the
final black hole state resulting from gravitational collapse can lose all information
about the collapsing body except its mass, angular momentum, and charge.

